\documentclass[12pt]{article}
\usepackage[a4paper, total={6.5in, 9in}]{geometry}
\usepackage{ebproof}
\usepackage{mathtools}
\usepackage{float}
\usepackage{amssymb}
\usepackage[dvipsnames]{xcolor}
\usepackage[T1]{fontenc}
\usepackage{colortbl}


\newcommand{\tx}{\text{\rm x}}
\newcommand{\ty}{\text{\rm y}}
\newcommand{\Beh}{\ensuremath{\mathcal{B}}}
\newcommand{\BX}{\ensuremath{\mathcal{X}}}
\newcommand{\Next}{\ensuremath{\mathcal{N}}}
\newcommand{\emptyenv}{\varnothing}

\newcommand{\Gammatop}{\ensuremath{\Xi}}
\newcommand{\Gammabot}{\ensuremath{\Gamma}}

\newcommand{\dom}{\texttt{dom}}
\newcommand{\locals}{\texttt{locals}}
\newcommand{\name}{\texttt{name}}
\newcommand{\atom}{\texttt{atom}}
\renewcommand{\mod}{\texttt{mod}}
\newcommand{\fun}{\texttt{fun}}
\newcommand{\Elixir}{\texttt{Elixir}}
\newcommand{\Self}{\texttt{Self}}%\_\_MODULE\_\_}}
\newcommand{\private}{\texttt{private\;}}
\newcommand{\type}{\texttt{type\;}}
\newcommand{\opaque}{\texttt{opaque\;}}
\newcommand{\wrap}{\texttt{wrap}}

\NewDocumentCommand{\Moddecl}{O{\tx} O{\overline {\mathcal{D}}} O{\BX} O{\tx'} O{t} }{
   \Lambda (\overline{#1 : \star}) . \left\{ #2 \right\}_{\overline{#3 [\overline{#4 = #5}]}}
}

\ebproofset{right label template=\small\textsc{\inserttext}}

%%%%%%%%%%%%%%%%%%%%%%%%%%%%%%%%%%%%%%%%%%%%%%%%%%%%%%%%%%%%%%%%%%%%%%%%%%%%%%%%
% Comments
%%%%%%%%%%%%%%%%%%%%%%%%%%%%%%%%%%%%%%%%%%%%%%%%%%%%%%%%%%%%%%%%%%%%%%%%%%%%%%%%
\definecolor{darkgreen}{rgb}{0, 0.2, 0.13}
\definecolor{lessdarkgreen}{rgb}{0, 0.3, 0.13}
\definecolor{bordeaux}{rgb}{0.5, 0, 0.13}
\definecolor{darkblue}{RGB}{51, 102, 255}
\definecolor{darkgray}{RGB}{99,102,106}

\newcommand{\beppe}[1]{{\color{bordeaux}\textbf{[ Beppe: } #1\textbf{]}}}
\newcommand{\aghilas}[1]{{\color{lessdarkgreen}\textbf{[ Aghilas: } #1\textbf{]}}}
\newcommand{\antonin}[1]{{\color{darkblue}\textbf{[ Antonin: } #1\textbf{]}}}
\newcommand{\deprecated}[1]{{\color{darkgray}\textbf{[ Deprecated: } #1\textbf{]}}}

\begin{document}

\section{Outline}

We propose a new type system for Elixir's modules.
The syntax language is available in file \texttt{melixir.pdf} on page 10.
We will recall what we think is necessary in this document.
A clean, self-sufficient document will be written later.

\paragraph{Elixir program}

We represent an Elixir program as a list of behaviours $\overline{\mathcal{B}}$ followed by a
list of bindings $\overline{B}$.

Checking a program involves checking that all behaviours and bindings are well-formed.
Informally, a behaviour is well-formed if all type and behaviour names are \textit{bound},
that is, they are declared before use. A binding is wellformed if all type, behaviour and module
names are bound and the values it defines are well typed.

\paragraph{Scoping rules}

We impose a significant limitation compared to current Elixir's functionality.
Elixir allows mutually recursive module definitions, but the system we propose
currently does not. As a result, at any point of the program, only preceding behaviours
and declarations can be used. In particular, because we separate behaviours and bindings
in two lists and we process behaviours first, behaviours cannot refer to modules at all.

Supporting mutually recursive modules would add significant complexity and is kept as
a work for later.

\paragraph{Environments}

We now define the type system formally below.

Recall the definition of behaviour, module and local environments:

\begin{center}
    \begin{tabular}{rl}
        $\Sigma ::=$    & $\overline{\mathcal B}$   \\[3mm]
        $\Gammatop ::=$ & $\overline{X : \Moddecl}$ \\[3mm]
        $\Gammabot ::=$ & $\overline{D}$
    \end{tabular}
\end{center}

where $D$ is a local declaration:

\begin{center}
    \begin{tabular}{rcll}
        $D$ & ::= & $ X : t$                   & \color{Gray} function parameter with module or behaviour type \\
            & |   & $x : \bigcap \overline{T}$ & \color{Gray} exports a value $x$ of type $\bigcap\overline T$ \\
            & |   & $\tx : \star$              & \color{Gray} exports an opaque type $x$                       \\
            & |   & $\tx : [=t]$               & \color{Gray} exports a transparent type $x = t$               \\
    \end{tabular}
\end{center}

and $\mathcal{D}$ is either a local declaration or a submodule declaration:

\begin{center}
    \begin{tabular}{rcl}
        $\mathcal{D}$ & ::= & $D$            \\
                      & |   & $X : \Moddecl$
    \end{tabular}
\end{center}

Empty lists and environments are denoted $\emptyenv$. We consider behaviour, module and local
environments as partial functions over $\BX$, $\Next$ and $\Next$ respectively,
where $\Next ::= N$ | $\tx$ ranges over declaration names.

By extension, we consider global environments as partial functions over $P$:

$$
    (\overline{\mathcal{D}})(X.P) = (\overline{\mathcal{D}})(X)(P)
$$


\paragraph{Submodules}

Recall module paths: $P ::= X$ | $P.X$. 

A module can be defined inside another module.
The absolute path of the inner module is obtained by prefixing
its name with the absolute path of the enclosing module.

To respect the scoping rules described above, the module environment must
reproduce this hierarchy and put inner module declarations inside the enclosing
module.
While necessary, this representation is unfortunate because it
may suggest that inner modules are members of their enclosing module.
This is not the case: the hierarchy is purely lexical.

Submodules do not inherit the local environment of their enclosing module.
Submodules declarations can never happen in types $\{\overline{D}\}$.
Similarly, $X : t$ can only be a function parameter; it is not a declaration that can be
exported. As a result, neither can it appear in types $\{\overline{D}\}$.

Admittedly, Elixir does allow references to a submodule from within its enclosing
module using relative paths.
However, this is only syntactic sugar. Relative paths are expanded to
their corresponding absolute path before they reach the type system.

\paragraph{Judgments}

We define the following judgements:

\begin{figure}[H]
    \begin{center}
        \begin{tabular}{lrl}
            Wellformedness of behaviour environments: &                                     & $\vdash \Sigma$                                             \\[2mm]
            Wellformedness of module environments:    & $\Sigma$                            & $\vdash \Gammatop$                                          \\[2mm]
            Wellformedness of local environments:     & $\Sigma; \Gammatop$                 & $\vdash \Gammabot$                                          \\[2mm]
            Wellformedness of types:                  & $\Sigma; \Gammatop; \Gammabot$      & $\vdash T$                                                  \\[2mm]
            Subtyping:                                & $\Sigma; \Gammatop; \Gammabot$      & $\vdash \bigcap \overline{T} \preceq \bigcap \overline{T'}$ \\[2mm]
            Expression typing:                        & $\Sigma; \Gammatop; \Gammabot$      & $\vdash E : \bigcap \overline{T} $                          \\[2mm]
            Binding typing:                           & $\Pi; \Sigma; \Gammatop; \Gammabot$ & $\vdash \overline{B} : \overline{\mathcal{D}} $             \\[2mm]
        \end{tabular}
    \end{center}
\end{figure}
where $\Pi$ is a possibly empty path:  $\Pi$ ::= $\epsilon$ | $P$.

The wellformedness of module environments will also need two auxiliary judgements, which will be defined in the corresponding section.

Let $\overline{\mathcal B}, \overline{B}$ be a program. The program passes type-checking if there exists a wellformed module environment $\Gammatop$ such that both
$$\vdash \overline{\mathcal B} \quad \text{ and } \quad \epsilon; \overline{\mathcal{B}}; \emptyenv; \emptyenv \vdash \overline{B} : \Gammatop$$ hold.

Note that we write $\overline{B} : \Gammatop$ because top-level bindings can only be module declarations.
More details about why we need to return it in the judgment are available in the section about Bindings.

\paragraph{Shadowing and alpha-renaming}

The type system strictly prohibits name shadowing.

This restriction is a direct consequence of representing
environments as partial functions that resolve to the right-most binding. Allowing shadowing would implicitly alter the meaning of any prior declarations that refer to the shadowed name.

While this correctly rejects modules that export the same name twice
and functions with duplicate parameter names, 
it can also be overly restrictive. It prevents function parameters
from shadowing local module declarations.
More critically, it makes it difficult to reference an external module
that shares declaration names with the current referencing module.

For now, we consider it is the responsibility of the programmer to
ensure names do not repeat. Later, we will formally define where and how
alpha-renaming can be applied to lift these unintended restrictions.


\section{Wellformedness}

In this section we give inference rules for the wellformedness judgements.

\subsection{Behaviour environments}

\[
    \begin{prooftree}
        \infer0{ \vdash \emptyenv }
    \end{prooftree}
    \qquad
    \begin{prooftree}
        \hypo{ \BX \notin \dom(\Sigma) }
        \hypo{ \Sigma; \emptyenv; \overline{\tx : \star} \vdash \{ \overline{D} \}}
        \infer2{ \vdash \Sigma, \BX : \Lambda (\overline{\tx : \star}).\{\overline{D}\} }
    \end{prooftree}
\]

$ \Sigma; \emptyenv; \overline{\tx : \star} \vdash \{ \overline{D} \}$ refers to the wellformedness of types.

\subsection{Module environments}

\[
    \begin{prooftree}
        \hypo { \epsilon; \Sigma; \emptyenv; \emptyenv \vdash \Gammatop }
        \infer1 { \Sigma \vdash \Gammatop }
    \end{prooftree}
\]

The hypothesis refers to the wellformedness of declaration lists, which is defined below. \\

Before we can define wellformedness of declaration lists, we need to introduce two new notations:
$\locals$ takes the list of declarations of a module and filters out submodule
declarations, so as to keep only local declarations. $(\overline{\mathcal{D}})[\Pi \texttt{+=} \mathcal{D}']$ adds a declaration
to a module in the module environment. The notations are formally defined below.

\begin{figure}[H]
    \begin{center}
        \begin{tabular}{rl}
            $\locals(\emptyenv)$                                        & $= \emptyenv$                          \\
            $\locals(D, \overline{\mathcal{D}})$                        & $= D, \locals(\overline{\mathcal{D}})$ \\
            $\locals\left(X : \Moddecl, \overline{\mathcal{D'}}\right)$ & $= \locals(\overline{\mathcal{D'}})$
        \end{tabular}
    \end{center}
    \caption{Meta-function $\locals(\overline{\mathcal{D}})$}
\end{figure}

\begin{figure}[H]
    \begin{center}
        \begin{tabular}{ll}$
            \left(\overline{\mathcal{D}}\right)\left[\epsilon \texttt{+=} X : \Moddecl[\tx][\overline{\mathcal{D}'}]\right]$ & = $\overline{\mathcal{D}}, X : \Moddecl[\tx][\overline{\mathcal{D}'}]$                              \\
            $\left(\overline{\mathcal{D}}, X : \Moddecl[\tx][\overline{\mathcal{D}'}]\right)[X \texttt{+=} \mathcal{D}'']$   & = $\overline{\mathcal{D}}, X : \Moddecl[\tx][\overline{\mathcal{D}'}, \mathcal{D}'']$               \\
            $\left(\overline{\mathcal{D}}, X : \Moddecl\right)[X.P \texttt{+=} \mathcal{D}']$                                & = $\overline{\mathcal{D}}, X : \Moddecl[\tx][(\overline{\mathcal{D}})[P \texttt{+=} \mathcal{D}']]$ \\
            undefined otherwise
        \end{tabular}
    \end{center}
    \caption{Notation $(\overline{\mathcal{D}})[\Pi \texttt{+=} \mathcal{D}']$}
\end{figure}

Note that we can only add a declaration to the last module in the list.

\subsection{Auxiliary judgements}

We define two new wellformedness judgements:
\begin{figure}[H]
    \begin{center}
        \begin{tabular}{lrl}
            Wellformedness of declaration lists    & $\Pi; \Sigma; \Gammatop; \Gammabot$ & $\vdash \overline{\mathcal{D}}$ \\[2mm]
            Wellformedness of module declarations: & $\Pi; \Sigma; \Gammatop$            & $\vdash X : \Moddecl$           \\
        \end{tabular}
    \end{center}
\end{figure}

The first judgement is not to be confused with the wellformedness of the type $\{\overline{D}\}$, which can only contain local declarations.

Note that there is no ambiguity between the wellformedness of a declaration list
that consists of a single module declaration $\left(\Pi; \Sigma; \Gammatop; \Gammabot \vdash X : \Moddecl\right)$
and the wellformedness of the module itself $\left(\Pi; \Sigma; \Gammatop \vdash X : \Moddecl\right)$
because the latter does not use the local environment $\Gammabot$.


\paragraph{Empty list}
\[
    \begin{prooftree}
        \hypo { \Sigma; \Gammatop \vdash \Gammabot }
        \infer1 { \Pi; \Sigma; \Gammatop; \Gammabot \vdash \emptyenv }
    \end{prooftree}
\]

\paragraph{Local declarations}

Intuitively, a local declaration is wellformed if we can add it to the
local environment built so far.

\[
    \begin{prooftree}
        \hypo { \Sigma; \Gammatop \vdash \Gammabot, D }
        \infer1 { P; \Sigma; \Gammatop; \Gammabot \vdash D }
    \end{prooftree}
\]

\[
    \begin{prooftree}
        \hypo { P; \Sigma; \Gammatop[P \texttt{+=} D]; \Gammabot, D \vdash \overline{\mathcal{D}'} }
        \infer1[$\overline{\mathcal{D}} \neq \emptyenv$] { P; \Sigma; \Gammatop; \Gammabot \vdash D, \overline{\mathcal{D}'} }
    \end{prooftree}
\]

Note that in the rules above, $P$ cannot be the empty path, because local declarations
can only appear inside a module. The special rule when there is a single local declaration is necessary to break
cycles in the definition.
Indeed, checking the wellformedness of the local environment will in turn
require checking the wellformedness of the module environment built so far.
We need this environment to be stricly smaller than the initial environment, otherwise we would get stuck.

\paragraph{Module declarations}

\[
    \begin{prooftree}
        \hypo { \Sigma; \Gammatop \vdash \Gammabot }
        \hypo { \Pi; \Sigma; \Gammatop \vdash X : \Moddecl }
        \infer2{ \Pi; \Sigma; \Gammatop; \Gammabot \vdash X : \Moddecl }
    \end{prooftree}
\]
\[
    \begin{prooftree}
        \hypo { \Pi; \Sigma; \Gammatop \vdash X : \Moddecl }
        \hypo { \Pi; \Sigma; \Gammatop\left[\Pi \texttt{+=} X : \Moddecl\right]; \Gammabot \vdash \overline{\mathcal{D}'} }
        \infer2 [$\overline{\mathcal{D}'} \neq \emptyenv$]{ \Pi; \Sigma; \Gammatop; \Gammabot \vdash X : \Moddecl, \overline{\mathcal{D}'} }
    \end{prooftree}
\]

By abuse of notation, we consider that $\epsilon.X$ denotes $X$. This avoids duplicating the rules.

\[
    \begin{prooftree}
        \hypo {  \Sigma; \Gammatop \vdash \overline{\tx : \star} }
        \hypo { \Pi.X \notin \dom(\Gammatop) }
        \infer2 { \Pi; \Sigma; \Gammatop \vdash X : \Lambda (\overline{\tx : \star}) . \{\} }
    \end{prooftree}
\]

\[
    \begin{prooftree}
        \hypo { \Pi.X; \Sigma; \Gammatop\left[\Pi \texttt{+=} X : \Lambda (\overline{\tx : \star}). \{\}\right]; \overline{\tx : \star} \vdash \overline{\mathcal{D} } }
        \hypo { \Sigma; \Gammatop\left[\Pi \texttt{+=} X : \Lambda (\overline{\tx : \star}). \left\{ \overline{\mathcal{D}}\right\}\right]; \overline{\tx : \star} \vdash \{ \locals( \overline{\mathcal{D}})\}_{\overline{\BX[\overline{\tx' = t}]}} }
        \infer2 { \Pi; \Sigma; \Gammatop \vdash X : \Moddecl }
    \end{prooftree}
\]

\antonin{

    Example of derivation:
    \[
        \begin{prooftree}
            \infer0{ \vdash \emptyenv }
            \infer1{ \emptyenv \vdash \emptyenv }
            \infer1[$A \notin \dom(\emptyenv)$]{ \epsilon; \emptyenv; \emptyenv \vdash A : \{ \} }
            \infer1[$\emptyenv; \emptyenv \vdash \emptyenv$]{  \epsilon; \emptyenv; \emptyenv; \emptyenv \vdash A : \{\} }
            \infer1{  \emptyenv \vdash A : \{\} }
            \infer1[$\tx \notin \dom(\emptyenv)$]{ \emptyenv; A : \{\} \vdash \tx : \star }
            \infer1{ A; \emptyenv; A : \{\}; \emptyenv \vdash \tx : \star }
            \infer1{ \epsilon; \emptyenv; \emptyenv; \emptyenv \vdash A : \{ \tx : \star \} }
            \infer1{ \emptyenv \vdash A : \{ \tx : \star \} }
            \infer1{ A; \emptyenv; A : \{ \tx : \star \} \vdash B : \{\} }
            \infer1{ A; \emptyenv; A : \{ \tx : \star \}; \tx : \star \vdash B : \{\} }
            \infer1{ A; \emptyenv; A : \{ \}; \emptyenv \vdash \tx : \star, B : \{\} }
            \infer1{ \epsilon; \emptyenv; \emptyenv \vdash A : \{ \tx : \star, B : \{\} \} }
            \infer1[$\emptyenv; \emptyenv \vdash \emptyenv$]{ \epsilon; \emptyenv; \emptyenv; \emptyenv \vdash A : \{ \tx : \star, B : \{\} \} }
            \infer1{ \emptyenv \vdash A : \{ \tx : \star, B : \{\} \} }
            \infer1{ \emptyenv; A : \{ \tx : \star, B : \{\}\} \vdash \emptyenv  }
            \infer1{ \emptyenv; A : \{ \tx : \star, B : \{\}\}; \emptyenv \vdash \langle A \rangle }
            \infer1{ \emptyenv; A : \{ \tx : \star, B : \{\}\}; \emptyenv \vdash A : \langle A \rangle }

            \infer0{ \emptyenv; A : \{ \tx : \star, B : \{\} \}; \emptyenv \vdash \langle A \rangle }
            \infer1{ \emptyenv; A : \{ \tx : \star, B : \{\} \}; \emptyenv \vdash \langle A \rangle \preceq \{ \tx : \star \} }
            \infer2{ \emptyenv; A : \{ \tx : \star, B : \{\} \}; \emptyenv \vdash A : \{ \tx : \star \} }
            \infer1 { A.B; \emptyenv; A : \{ \tx : \star, B : \{\} \}; \emptyenv \vdash  A.\tx }
            \infer1[$\tx \notin \dom(\emptyenv)$] { A.B; \emptyenv; A : \{ \tx : \star, B : \{\} \} \vdash  \ty : [= A.\tx] }
            \infer1[$\emptyenv; \emptyenv \vdash \emptyenv$]{ A.B; \emptyenv; A : \{ \tx : \star, B : \{\} \}; \emptyenv \vdash  \ty : [= A.\tx] }
            \infer1 { A; \emptyenv; A : \{ \tx : \star \} \vdash   B : \{ \ty : [= A.\tx] \}}

            \infer0 { (*) }

            \infer2 { A; \emptyenv; A : \{ \tx : \star \}; \tx : \star  \vdash   B : \{ \ty : [= A.\tx] \}, \tx' : [= \tx] }
            \infer1 { A; \emptyenv; A : \{\}; \emptyenv  \vdash  \tx : \star, B : \{ \ty : [= A.\tx] \}, \tx' : [= \tx] }            \infer1 { \epsilon; \emptyenv; \emptyenv \vdash A : \{ \tx : \star, B : \{ \ty : [= A.\tx] \}, \tx' : [= \tx] \} }

            \infer1[$\emptyenv; \emptyenv \vdash \emptyenv$] { \epsilon; \emptyenv; \emptyenv; \emptyenv \vdash A : \{ \tx : \star, B : \{ \ty : [= A.\tx] \}, \tx' : [= \tx] \} }
            \infer1 { \emptyenv \vdash A : \{ \tx : \star, B : \{ \ty : [= A.\tx] \}, \tx' : [= \tx] \} }
        \end{prooftree}
    \] \\

    \[
        \begin{prooftree}
            \infer0{ A; \emptyenv; A: \{ \tx : \star\} \vdash  B : \{ \ty : [= A.x] \} }

            \infer0{ \emptyenv \vdash A : \{ \tx : \star \} }
            \infer1{ \emptyenv; A : \{ \tx : \star \} \vdash \emptyenv }
            \infer1[$\tx \notin \dom(\emptyenv)$]{ \emptyenv; A: \{ \tx : \star\} \vdash  \tx : \star }

            \infer2{ A; \emptyenv; A: \{ \tx : \star\}; \tx : \star \vdash  B : \{ \ty : [= A.x] \} }
            \infer1{ A; \emptyenv; A: \{\}; \emptyenv \vdash \tx : \star, B : \{ \ty : [= A.x] \} }
            \infer1{ \epsilon; \emptyenv; \emptyenv; \emptyenv \vdash A : \{ \tx : \star, B : \{ \ty : [= A.x] \}\} }
            \infer1{ \emptyenv \vdash A : \{ \tx : \star, B : \{ \ty : [= A.x] \}\} }
            \infer1{ \emptyenv; A : \{ \tx : \star, B : \{ \ty : [= A.x] \}\} \vdash \emptyenv }
            \infer1[$\tx \notin \dom(\emptyenv)$]{  \emptyenv; A : \{ \tx : \star, B : \{ \ty : [= A.x] \}\} \vdash \tx : \star }
            \infer1[$\tx \in \dom(\tx : \star)$]{  \emptyenv; A : \{ \tx : \star, B : \{ \ty : [= A.x] \}\}; \tx : \star \vdash \tx }
            \infer1[$\tx' \notin \dom(\tx : \star)$]{  \emptyenv; A : \{ \tx : \star, B : \{ \ty : [= A.x] \}\} \vdash \tx : \star, \tx' : [= \tx] }
            \infer1{ (*) = A; \emptyenv; A : \{ \tx : \star, B : \{ \ty : [= A.x] \}\}; \tx : \star  \vdash  \tx' : [= \tx]}
        \end{prooftree}
    \]

}

\subsection{Local type environments}

\[
    \begin{prooftree}
        \hypo{ \Sigma \vdash \Gammatop}
        \infer1{ \Sigma; \Gammatop \vdash \emptyenv }
    \end{prooftree}
\]

\[
    \begin{prooftree}
        \hypo{ \Sigma; \Gammatop; \Gammabot \vdash t \preceq \mod }
        \hypo{ X \notin \dom(\Gammabot) }
        \infer2{ \Sigma; \Gammatop \vdash \Gammabot, X : t }
    \end{prooftree}
\]
\[
    \begin{prooftree}
        \hypo{ \Sigma; \Gammatop \vdash \Gammabot }
        \hypo{ \tx \notin \dom(\Gammabot) }
        \infer2{ \Sigma; \Gammatop \vdash \Gammabot, \tx : \star }
    \end{prooftree}
\]
\[
    \begin{prooftree}
        \hypo{ \Sigma; \Gammatop; \Gammabot \vdash t }
        \hypo{ \tx \notin \dom(\Gammabot) }
        \infer2{ \Sigma; \Gammatop \vdash \Gammabot, \tx : [= t] }
    \end{prooftree}
\]
\[
    \begin{prooftree}
        \hypo{ \overline{\Pi; \Sigma; \Gammatop; \Gammabot \vdash T} }
        \hypo{ x \notin \dom(\Gammabot) }
        \infer2[$\overline{T} \neq \emptyset$]{ \Pi; \Sigma; \Gammatop \vdash \Gammabot, x : \bigcap \overline{T} }
    \end{prooftree}
\] \\

Note that (if we defined inference rules properly),
$\Sigma; \Gammatop; \Gammabot \vdash t \preceq t'$ implies $\Sigma; \Gammatop; \Gammabot \vdash t$,
which in turn implies $\Sigma; \Gammatop \vdash \Gammabot$;
that is why we did not mention this as an explicit condition in the rules above.


\subsection{Pretypes}

Recall instantiated module paths: $M ::= X$ | $\Elixir.P[\overline{\tx = t}]$.
This may come as a surprise because we mentioned in the beginning of the document
that modules can only be accessed using their absolute path.
However, here, $X$ is not a relative path to a submodule; this is the name of a function parameter.
As such, it cannot be instantiated again because parameters must already have a concrete, fully
instantiated type (parametric types are not first-class).
Likewise, submodules are not included in the parent module's signature, so $X.P$ (where $X$ is a
function parameter) would not make sense.

Absolute paths are qualified with the $\Elixir$ prefix to remove ambiguity between the function parameter $X$
and the top-level module with name $X$.

\[
    \begin{prooftree}
        \hypo { \Sigma; \Gammatop \vdash \Gammabot }
        \infer1{ \Sigma; \Gammatop; \Gammabot \vdash \tau }
    \end{prooftree}
    \qquad
    \begin{prooftree}
        \hypo{ \overline{ \Sigma; \Gammatop; \Gammabot \vdash t}}
        \hypo{ \Sigma; \Gammatop; \Gammabot \vdash t'}
        \infer2{ \Sigma; \Gammatop; \Gammabot \vdash (\overline{t}) \rightarrow t' }
    \end{prooftree}
    \qquad
    \begin{prooftree}
        \hypo{ \Sigma; \Gammatop \vdash \Gammabot }
        \hypo{ \overline{ \Sigma; \Gammatop; \Gammabot \vdash t} }
        \infer2{ \Sigma; \Gammatop; \Gammabot \vdash  \%\{ \overline{k: t} \} }
    \end{prooftree}
\]

Keys $\overline{k}$ also need to be pairwise distinct. \\

\[
    \begin{prooftree}
        \hypo{ \Sigma; \Gammatop; \Gammabot \vdash t }
        \hypo{ \Sigma; \Gammatop; \Gammabot \vdash t' }
        \infer2{ \Sigma; \Gammatop; \Gammabot \vdash  t \vee  t' }
    \end{prooftree}
    \qquad
    \begin{prooftree}
        \hypo{ \Sigma; \Gammatop; \Gammabot \vdash t }
        \infer1{ \Sigma; \Gammatop; \Gammabot \vdash  \lnot t }
    \end{prooftree}
\]

\[
    \begin{prooftree}
        \hypo{  \Sigma; \Gammatop \vdash \Gammabot }
        \hypo{ \tx \in \dom(\Gammabot) }
        \infer2{ \Sigma; \Gammatop; \Gammabot \vdash \tx }
    \end{prooftree}
    \qquad
    \begin{prooftree}
        \hypo {  \Sigma; \Gammatop \vdash \Gammabot }
        \hypo { \Sigma(\BX) = \BX : \Lambda \left( \overline{\tx : \star} \right) \{\overline D \}}
        \hypo { \overline { \Sigma; \Gammatop; \Gammabot \vdash t} }
        \infer3{  \Sigma; \Gammatop; \Gammabot \vdash \BX [\overline{\tx = t}]  }
    \end{prooftree}
\]

\[
    \begin{prooftree}
        \hypo{ \Sigma; \Gammatop; \Gammabot \vdash M : \{ \tx : \star \} }
        \infer1{ \Sigma; \Gammatop; \Gammabot \vdash M.\tx }
    \end{prooftree}
\]

In the above rule, $\Sigma; \Gammatop; \Gammabot \vdash M : \{ \tx : \star \}$ is a typing judgement.
Thanks to subtyping, $\{ \tx : \star \}$ may not be exactly the type associated with $M$ in $\Gammatop$ or $\Gammabot$.


\subsection{Singleton module type}

\[
    \begin{prooftree}
        \hypo { \Sigma; \Gammatop \vdash \Gammabot }
        \hypo { \Gammatop(P) = X : \Moddecl[\tx] [\overline{\mathcal{D}}][\BX][\tx'][t']}
        \hypo { \overline{\Sigma; \Gammatop; \Gammabot \vdash t} }
        \infer3{ \Sigma; \Gammatop; \Gammabot \vdash \langle \Elixir.P[\overline{\tx = t}] \rangle }
    \end{prooftree}
\]

\[
    \begin{prooftree}
        \hypo {\Gammabot(X) = X : t }
        \hypo { \Sigma; \Gammatop; \Gammabot \vdash t \preceq \mod }
        \infer2{ \Sigma; \Gammatop; \Gammabot \vdash \langle X \rangle }
    \end{prooftree}
\]

\subsection{Function on modules}

\[
    \begin{prooftree}
        \hypo { \Sigma; \Gammatop; \Gammabot, \overline{N : T} \vdash  T' }
        \infer1{ \Sigma; \Gammatop; \Gammabot  \vdash (\overline{N : T}) \rightarrow T' }
    \end{prooftree}
\]

Note that $\Sigma; \Gammatop; \Gammabot, \overline{N : T} \vdash  T'$
implies that $\Gammabot, \overline{N : T}$ is wellformed, which in turn implies
that parameters $\overline{N : T}$ are wellformed and the $\overline{N}$ are pairwise distinct.

\subsection{List of declarations}

The name "declaration list" that refers to the type
$\{ \overline{D} \}$ is not to be confused with the declaration list
$\overline{\mathcal{D}}$ which appears in the module environment.

We define the function $\name(D)$ as $\name: \left\{\begin{array}{ll}
        X : t                    & \mapsto X   \\
        x : \bigcap \overline{T} & \mapsto x   \\
        \tx : \star              & \mapsto \tx \\
        \tx : [=t]               & \mapsto \tx
    \end{array}\right.$  \\[3mm]

\[
    \begin{prooftree}
        \hypo{ \Sigma; \Gammatop \vdash \Gammabot }
        \infer1{ \Sigma; \Gammatop; \Gammabot \vdash \{ \} }
    \end{prooftree}
\]
\[
    \begin{prooftree}
        \hypo{ \Sigma; \Gammatop; \Gammabot, D \vdash \{ \overline{D'} \} }
        \infer1[\;($ D \neq X : t$)]{ \Sigma; \Gammatop; \Gammabot \vdash \{ D, \overline{D'} \} }
    \end{prooftree}
\] \\

\[
    \begin{prooftree}
        \hypo{ \Sigma; \Gammatop; \Gammabot \vdash \{\overline{D} \}}
        \hypo{\overline{\Sigma; \Gammatop; \Gammabot \vdash \BX[\overline{\tx = t}]} }
        \hypo{ \forall i \neq j. \; \overline{\BX}^i \# \overline{\BX}^j }
        \hypo{\overline{\Sigma; \Gammatop; \Gammabot \vdash \{\overline{D}\} \preceq  \Sigma(\BX)[\overline{\tx := t}]}}
        \infer4{ \Sigma; \Gammatop; \Gammabot \vdash  \{ \overline{D} \}_{\overline{\BX[\overline{\tx = t}]}}}
    \end{prooftree}
\] \\


where, if we denote $\Sigma(\BX) = \BX : \Lambda (\overline{\tx : \star}) \{\overline{D}\}$ and
$\Sigma(\BX') = \BX' : \Lambda (\overline{\tx' : \star}) \{\overline{D'}\}$:
\begin{itemize}
    \item $\BX \# \BX'$ if and only if $\BX \neq \BX'$ and $\dom(\overline D) \cap \dom(\overline{D'}) = \emptyset$
    \item $\Sigma(\BX)[\overline{\tx := t}]$ is capture-avoiding substitution
\end{itemize}

The substitution is to be defined more formally. In particular, it will not be defined
if one of the $t$ contains a type name that is also in $\dom(\overline{D})$.

For now, we assume names never collide.
Also, note that the substitution does not traverse the dot operator (that is, $M.\tx[t/\tx]$ is still $M.\tx$).

\section{Subtyping}

Elixir's semantic subtyping will be denoted $ t \vartriangleleft t'$.

Recall the syntactic subtyping judgement: $\Sigma; \Gammatop; \Gammabot \vdash \bigcap \overline{T} \preceq \bigcap \overline{T'}$.
We will denote $\bigcap \overline{T} \simeq \bigcap \overline{T'}$ when both $\bigcap \overline{T} \preceq \bigcap \overline{T'}$ and $\bigcap \overline{T'} \preceq \bigcap \overline{T}$ hold.

\subsection{General subtyping rules}

\[
    \begin{prooftree}
        \hypo{ \Sigma; \Gammatop \vdash \Gammabot }
        \hypo{ t \vartriangleleft t' }
        \infer2{ \Sigma; \Gammatop; \Gammabot \vdash t \preceq t' }
    \end{prooftree}
\] \\

\[
    \begin{prooftree}
        \hypo{ \Sigma; \Gammatop; \Gammabot \vdash \{ \overline{D} \}_{\overline{\BX[\overline{\tx = t}]} } }
        \hypo{ \Sigma; \Gammatop; \Gammabot \vdash \{ \overline{D'} \}_{\overline{\BX[\overline{\tx' = t'}]} } }
        \hypo{ \Sigma; \Gammatop; \Gammabot \vdash \{ \overline{D} \} \cap \bigcap \overline{\BX[\overline{\tx = t}]} \preceq \{ \overline{D'} \} \cap \bigcap \overline{\BX'[\overline{\tx' = t'}]}}
        \infer3{ \Sigma; \Gammatop; \Gammabot \vdash \{ \overline{D} \}_{\overline{\BX[\overline{\tx = t}]} } \preceq \{ \overline{D'} \}_{\overline{\BX'[\overline{\tx' = t'}]} } }
    \end{prooftree}
\] \\

\[
    \begin{prooftree}
        \hypo{ \Sigma; \Gammatop; \Gammabot \vdash \{ \overline{D} \}_{\overline{\BX[\overline{\tx = t}]} }}
        \infer1{ \Sigma; \Gammatop; \Gammabot \vdash \{ \overline{D} \}_{\overline{\BX[\overline{\tx = t}]} } \preceq \bigcap \overline{\BX[\overline{\tx = t}]} }
    \end{prooftree}
\] \\

\[
    \begin{prooftree}
        \hypo{ \Sigma; \Gammatop \vdash \Gammabot }
        \hypo{ \overline{\Sigma; \Gammatop; \Gammabot \vdash T} }
        \hypo{ \forall j. \exists i. \; \Sigma; \Gammatop; \Gammabot \vdash \overline{T}^i \preceq \overline{T'}^j }
        \infer3{ \Sigma; \Gammatop; \Gammabot \vdash \bigcap\limits_i \overline{T}^i \preceq \bigcap\limits_j \overline{T'}^j}
    \end{prooftree}
\] \\

We dont mention $\overline{\Sigma; \Gammatop; \Gammabot \vdash T'}$
as an explicit condition because $\overline{T'}^j$ ranges over all $\overline{T'}$.
As a result, the hypothesis $\overline{T}^i \preceq \overline{T'}^j$ already implies that
all the $\overline{T'}^j$ are wellformed.

Note that with this definition, the empty intersection is a supertype of all intersections,
and therefore is a supertype of all types. This is consistent with the convention that
the empty intersection is equivalent to the top type $\lnot \mathbb{O}$.

The other direction might (or might not) be useful later:

\[
    \begin{prooftree}
        \hypo{ \Sigma; \Gammatop; \Gammabot \vdash t \wedge t'}
        \infer1{ \Sigma; \Gammatop; \Gammabot \vdash t \cap t' \preceq t \wedge t' }
    \end{prooftree}
\] \\


\[
    \begin{prooftree}
        \hypo{ \Sigma; \Gammatop; \Gammabot \vdash (\overline{N : t }) \rightarrow t' }
        \hypo{ \Sigma; \Gammatop; \Gammabot \vdash (\overline{t}) \rightarrow t'}
        \infer2{ \Sigma; \Gammatop; \Gammabot \vdash (\overline{t}) \rightarrow t' \preceq (\overline{N : t }) \rightarrow t' }
    \end{prooftree}
\] \\

This rule is necessary to prove for instance:
\begin{align*}
    (\BX, \texttt{integer()}) \rightarrow \texttt{boolean()} & \preceq (N1: \BX, N2: \texttt{integer()}) \rightarrow \texttt{boolean()} \\
                                                             & \preceq (N1 : \BX, N2 : N1.\tx) \rightarrow N2.\tx'
\end{align*}

where $N1.\tx : [= \texttt{integer()}], N2.\tx' : [= \texttt{boolean()}]$ \\

\antonin{Do we need the other direction?

\[
    \begin{prooftree}
        \hypo{ \Sigma; \Gammatop; \Gammabot \vdash (\overline{N : t }) \rightarrow t' }
        \hypo{ \Sigma; \Gammatop; \Gammabot \vdash (\overline{t}) \rightarrow t'}
        \infer2{ \Sigma; \Gammatop; \Gammabot \vdash (\overline{N : t }) \rightarrow t' \preceq (\overline{t}) \rightarrow t'  }
    \end{prooftree}
\] \\

}

\[
    \begin{prooftree}
        \hypo { \forall i. \; \Sigma; \Gammatop; \Gammabot, \overline{N : T'}^{0..i} \vdash \overline{T'}^{i} \preceq  \overline{T}^{i} }
        \hypo{ \Sigma; \Gammatop; \Gammabot, \overline{N : T'} \vdash T_r \preceq T_r' }
        \infer2{ \Sigma; \Gammatop; \Gammabot \vdash (\overline{N : T}) \rightarrow T_r \preceq (\overline{N : T'}) \rightarrow T_r' }
    \end{prooftree}
\] \\

\antonin{We will probably need a rule to rename parameter names here, because it's
    not as trivial as it seems because of the risk of capture...}

\[
    \begin{prooftree}
        \hypo{ \Sigma; \Gammatop; \Gammabot \vdash \BX[\overline{\tx = t}]}
        \infer1{ \Sigma; \Gammatop; \Gammabot \vdash \BX[\overline{\tx = t}] \preceq \Sigma(\BX)[\overline{\tx := t}]}
    \end{prooftree}
\] \\


\[
    \begin{prooftree}
        \hypo{ \Sigma; \Gammatop; \Gammabot \vdash \{ \overline{D} \} }
        \infer1{ \Sigma; \Gammatop; \Gammabot \vdash \{ \overline{D} \} \preceq  \mod }
    \end{prooftree}
\] \\

\[
    \begin{prooftree}
        \hypo { \Sigma; \Gammatop; \Gammabot \vdash \langle \Elixir.P[\overline{\tx = t}] \rangle }
        \infer1{ \Sigma; \Gammatop; \Gammabot \vdash \langle \Elixir.P[\overline{\tx = t}] \rangle \preceq \Gammatop(P)[\overline{\tx := t}] }
    \end{prooftree}
\] \\

where $\left(X : \Moddecl[\tx][\overline{\mathcal{D}}][\BX][\tx'][t']\right)[\overline{\tx := t}] = \left(\left\{ \locals(\overline{\mathcal{D}}) \right\}_{\overline{\BX[\overline{\tx' = t'}]}}\right)[\overline{\tx := t}]$

\[
    \begin{prooftree}
        \hypo { \Sigma; \Gammatop \vdash \Gammabot }
        \hypo { \Gammabot(X) = X : t }
        \infer2{ \Sigma; \Gammatop; \Gammabot \vdash \langle X \rangle \preceq t}
    \end{prooftree}
\] \\

\subsection{Declaration lists}

The conditions $\tx \notin \dom(\overline{D})$ and $x \notin \dom(\overline{D})$ are there to ensure that the
enriched type $\{ \tx : ..., \overline{D} \}$ or $\{ x : ..., \overline{D} \}$ in the conclusion of the rule is still wellformed.


\paragraph{Base case}

\[
    \begin{prooftree}
        \hypo{ \Sigma; \Gammatop; \Gammabot \vdash \{ \overline{D} \} }
        \infer1{ \Sigma; \Gammatop; \Gammabot \vdash \{ \overline{D} \} \preceq \{ \} }
    \end{prooftree}
\]

\paragraph{Skip left declarations}

\[
    \begin{prooftree}
        \hypo{ \Sigma; \Gammatop; \Gammabot, \tx : \star  \vdash \{ \overline{D} \} \preceq \{\overline{D'}\} }
        \infer1{ \Sigma; \Gammatop; \Gammabot \vdash \{ \tx : \star, \overline{D} \} \preceq \{ \overline{D'} \} }
    \end{prooftree}
\]

\[
    \begin{prooftree}
        \hypo{ \Sigma; \Gammatop; \Gammabot \vdash \{ \overline{D}\}[\tx := t] \preceq \{\overline{D'}\} }
        \hypo{ \tx \notin \dom(\Gammabot) \cup \dom(\overline{D}) \cup \dom(\overline{D'}) }
        \infer2{ \Sigma; \Gammatop; \Gammabot \vdash \{ \tx : [= t], \overline{D} \} \preceq \{ \overline{D'} \} }
    \end{prooftree}
\]

\[
    \begin{prooftree}
        \hypo{ \Sigma; \Gammatop; \Gammabot \vdash \{ \overline{D} \} \preceq \{\overline{D'}\} }
        \hypo{ x \notin \dom(\Gammabot) \cup \dom(\overline{D}) \cup \dom(\overline{D'}) }
        \infer2{ \Sigma; \Gammatop; \Gammabot \vdash \{ x : \bigcap \overline{T}, \overline{D} \} \preceq \{ \overline{D'} \} }
    \end{prooftree}
\]

\paragraph{Induction}

\[
    \begin{prooftree}
        \hypo{ \Sigma; \Gammatop; \Gammabot, \tx : \star \vdash \{\overline{D}\} \preceq \{\overline{D'}\} }
        \infer1{ \Sigma; \Gammatop; \Gammabot \vdash \{ \tx : \star, \overline{D} \} \preceq \{ \tx : \star, \overline{D'} \} }
    \end{prooftree}
\] \\

\[
    \begin{prooftree}
        \hypo{ \Sigma; \Gammatop; \Gammabot, \tx : [= t] \vdash \{\overline{D}\} \preceq \{\overline{D'} \} }
        \infer1{ \Sigma; \Gammatop; \Gammabot \vdash \{ \tx : [= t], \overline{D} \} \preceq \{ \tx : \star, \overline{D'} \} }
    \end{prooftree}
\] \\

\[
    \begin{prooftree}
        \hypo{ \Sigma; \Gammatop; \Gammabot \vdash \{\overline{D}\}[\tx := t] \preceq \{\overline{D'}\}[\tx := t] }
        \hypo{ \tx \notin \dom(\Gammabot) \cup \dom(\overline{D}) \cup \dom(\overline{D'}) }
        \infer2{ \Sigma; \Gammatop; \Gammabot \vdash \{ \tx : [= t], \overline{D} \} \preceq \{ \tx : [= t'], \overline{D'} \} }
    \end{prooftree}
\] \\

We intentionnally do not enforce $t \preceq t'$ in the rule above because $\tx$ could be used
in a contravariant way.

\[
    \begin{prooftree}
        \hypo{ \Sigma; \Gammatop; \Gammabot \vdash \bigcap \overline{T} \preceq \bigcap \overline{T'} }
        \hypo{ \Sigma; \Gammatop; \Gammabot, x : \bigcap \overline{T} \vdash \{\overline{D}\} \preceq \{\overline{D'}\} }
        \infer2{ \Sigma; \Gammatop; \Gammabot \vdash \{ x: \bigcap \overline{T}, \overline{D} \} \preceq \{ x: \bigcap \overline{T'}, \overline{D'} \} }
    \end{prooftree}
\] \\

\paragraph{Reordering}

\[
    \begin{prooftree}
        \hypo{ \Sigma; \Gammatop; \Gammabot \vdash \{ D_1, D_2, \overline{D} \} \preceq \{ \overline{D'} \} }
        \hypo{ \Sigma; \Gammatop \vdash \Gammabot, D_2, D_1 }
        \infer2{ \Sigma; \Gammatop; \Gammabot \vdash \{ D_2, D_1, \overline{D} \} \preceq \{ \overline{D'} \} }
    \end{prooftree}
\] \\

\section{Expressions}

Recall the expression typing judgement: $\Sigma; \Gammatop; \Gammabot \vdash E : \bigcap \overline{T}$.

\[
    \begin{prooftree}
        \hypo{ \Sigma; \Gammatop; \Gammabot \vdash E : \bigcap \overline{T} }
        \hypo{ \Sigma; \Gammatop; \Gammabot \vdash \bigcap \overline{T}\preceq \bigcap \overline{T'} }
        \infer2{ \Sigma; \Gammatop; \Gammabot \vdash E: \bigcap \overline{T'}}
    \end{prooftree}
\] \\

\[
    \begin{prooftree}
        \hypo{ \Sigma; \Gammatop \vdash \Gammabot }
        \hypo{ \Gammabot(x) = x : \bigcap \overline{T} }
        \infer2{ \Sigma; \Gammatop; \Gammabot \vdash  x : \bigcap \overline{T} }
    \end{prooftree}
\] \\

\[
    \begin{prooftree}
        \hypo { \Sigma; \Gammatop; \Gammabot \vdash \%\{ \overline{k : t} \} }
        \hypo { \overline{\Sigma; \Gammatop; \Gammabot \vdash E : t} }
        \infer2{ \Sigma; \Gammatop; \Gammabot \vdash \%\{\overline{k = E}\} : \%\{\overline{k : t}\} }
    \end{prooftree}
\] \\

\[
    \begin{prooftree}
        \hypo { \overline{ \Sigma; \Gammatop; \Gammabot, \overline{x : t} \vdash E: t'} }
        \infer1[$\overline{t'} \neq \emptyset$]{ \Sigma; \Gammatop; \Gammabot \vdash \bigwedge \overline{(\overline{t}) \rightarrow t'} \text{ fn } \overline{x} \rightarrow E : \bigwedge \overline{ (\overline{t}) \rightarrow t'} }
    \end{prooftree}
\] \\

\[
    \begin{prooftree}
        \hypo { \overline{ \Sigma; \Gamma, \overline{N : T} \vdash E: T'} }
        \infer1[$\overline{T'} \neq \emptyset$]{ \Sigma; \Gammatop; \Gammabot \vdash \bigcap \overline{(\overline{N : T}) \rightarrow T'} \text{ fn } \overline{N} \rightarrow E : \bigcap \overline{(\overline{N : T }) \rightarrow T'} }
    \end{prooftree}
\] \\

\[
    \begin{prooftree}
        \hypo { \Sigma; \Gammatop; \Gammabot \vdash E : T }
        \hypo { \Sigma; \Gammatop; \Gammabot, N : T \vdash E' : T' }
        \infer2{ \Sigma; \Gammatop; \Gammabot \vdash \text{let } N = E \text{ in } E' : T' }
    \end{prooftree}
\] \\

\[
    \begin{prooftree}
        \hypo { \overline{\Sigma; \Gammatop; \Gammabot \vdash N : T[\overline{N/N'}]} }
        \hypo { \Sigma; \Gammatop; \Gammabot \vdash E : (\overline{N' : T}) \rightarrow T'}
        \infer2{ \Sigma; \Gammatop; \Gammabot \vdash E(\overline{N}) : T'[\overline{N/N'}] }
    \end{prooftree}
\] \\

Since neither forward references nor variable shadowing is allowed, $\overline{N'}^i$
cannot appear in $\overline{T}^j$ for $j < i$. As a result, $\overline{T[\overline{N/N'}]}^i$
is safe and equivalent to  $\overline{T[\overline{N/N'}^{0..i}]}^i$, which we think is less legible.


The hypothesis $\overline{\Sigma; \Gammatop; \Gammabot \vdash N : T[\overline{N/N'}]}$ means that
$\Sigma; \Gammatop; \Gammabot \vdash \overline{N}^i : \overline{T[\overline{N/N'}]}^i$ must hold for all $i$
without changing $\Gammabot$. It is not to be confused with the wellformedness judgement
$\Sigma; \Gammatop \vdash \Gammabot, \overline{N : T[\overline{N/N'}]}$.

We do not give a typing judgement for the application where the function has type $(\overline{t}) \rightarrow t'$
because it is a subtype of $(\overline{N : t}) \rightarrow t'$, as defined above.

\[
    \begin{prooftree}
        \hypo { \Sigma; \Gammatop; \Gammabot \vdash E : \%\{ k : t, .. \}  }
        \infer1{ \Sigma; \Gammatop; \Gammabot \vdash E.k : t }
    \end{prooftree}
\] \\

\[
    \begin{prooftree}
        \hypo { \Sigma; \Gammatop; \Gammabot \vdash E' : T }
        \hypo { \Sigma; \Gammatop; \Gammabot \vdash E'' : T  }
        \infer2{ \Sigma; \Gammatop; \Gammabot \vdash \left((E \in t) ? E' : E''\right) : T }
    \end{prooftree}
\] \\

We may also propose a better version where we give a name to the expression being tested:
\[
    \begin{prooftree}
        \hypo { \Sigma; \Gammatop; \Gammabot \vdash E_1 : t'}
        \hypo { \Sigma; \Gammatop; \Gammabot, N : t' \wedge t \vdash E_2 : T' }
        \hypo { \Sigma; \Gammatop; \Gammabot, N : t' \wedge \lnot t \vdash E_3 : T'  }
        \infer3{ \Sigma; \Gammatop; \Gammabot \vdash \left((N = E_1 \in t) ? E_2 : E_3\right) : T' }
    \end{prooftree}
\]
\[
    \begin{prooftree}
        \hypo { \Sigma; \Gammatop; \Gammabot \vdash M : \left\{ \overline{D}, x : \bigcap \overline{T} \right\} }
        \infer1{ \Sigma; \Gammatop; \Gammabot \vdash M.x : \bigcap \overline{\wrap_{M; \overline{D}}(T)} }
    \end{prooftree}
\] \\

where $\wrap$ qualifies with `$M.$` all the types declared in $\overline{D}$:
\begin{align*}
    \wrap_{M; \emptyenv}(T)                                                   & = T                                     \\
    \wrap_{M; \left(\tx \;:\; \star,\; \overline{D}\right)}(T)                & = \wrap_{M; \overline{D}}(T[M.\tx/\tx]) \\
    \wrap_{M; \left(\tx \;:\; [= t],\;  \overline{D}\right)}(T)               & = \wrap_{M; \overline{D}}(T[M.\tx/\tx]) \\
    \wrap_{M; \left(x \;:\; \bigcap \overline{T'},\;  \overline{D}\right)}(T) & = \wrap_{M; \overline{D}}(T)            \\
\end{align*}

Note that $\overline{D}$ appears inside a type $\{ \overline{D}, x : \bigcap \overline{T} \}$.
As a result, it cannot contain $X : t$; that is why we did not define $\wrap$ on $X : t$.

\[
    \begin{prooftree}
        \hypo { \Sigma; \Gammatop; \Gammabot \vdash \langle M \rangle }
        \infer1{ \Sigma; \Gammatop; \Gammabot \vdash M : \langle M \rangle }
    \end{prooftree}
\]

\[
    \begin{prooftree}
        \hypo{ \Sigma; \Gammatop \vdash \Gammabot }
        \hypo{ \Gammabot(X) = X : t }
        \infer2{ \Sigma; \Gammatop; \Gammabot \vdash X : t }
    \end{prooftree}
\]


\section{Bindings}

Recall the shape of the typing judgement for bindings: $\Pi; \Sigma; \Gammatop; \Gammabot \vdash \overline{B} : \overline{\mathcal{D}}$.

\[
    \begin{prooftree}
        \hypo{ \Sigma; \Gammatop \vdash \Gammabot }
        \infer1{ \Pi; \Sigma; \Gammatop; \Gammabot \vdash \emptyenv : \emptyenv}
    \end{prooftree}
\] \\

\[
    \begin{prooftree}
        \hypo{ \Sigma; \Gammatop; \Gammabot \vdash v: \bigcap \overline{T} }
        \hypo{ P; \Sigma; \Gammatop\left[P \texttt{+=} x : \bigcap \overline{T}\right]; \left(\Gammabot, x : \bigcap \overline{T}\right) \vdash \overline{B} : \overline{\mathcal{D}} }
        \infer2{ P; \Sigma; \Gammatop; \Gammabot \vdash \left(x = v, \overline{B}\right) : \left(x : \bigcap \overline{T}, \overline{\mathcal{D}}\right) }
    \end{prooftree}
\]  \\
\[
    \begin{prooftree}
        \hypo{ \Sigma; \Gammatop; \Gammabot \vdash v : \bigcap \overline{T}}
        \hypo{ P; \Sigma; \Gammatop; \left(\Gammabot, x : \bigcap \overline{T}\right) \vdash \overline{B} : \overline{\mathcal{D}}}
        \infer2{ P; \Sigma; \Gammatop; \Gammabot \vdash \left(\private x = v, \overline{B}\right): \overline{\mathcal{D}} }
    \end{prooftree}
\] \\
\[
    \begin{prooftree}
        \hypo{ P; \Sigma; \Gammatop\left[P \texttt{+=} \tx : [= t]\right]; \left(\Gammabot, \tx : [= t]\right) \vdash \overline{B} : \overline{\mathcal{D}}}
        \infer1{ P; \Sigma; \Gammatop; \Gammabot \vdash \left(\type \tx = t, \overline{B}\right) : \left(\tx : [= t], \overline{\mathcal{D}}\right) }
    \end{prooftree}
\] \\
\[
    \begin{prooftree}
        \hypo{ P; \Sigma; \Gammatop\left[P \texttt{+=} \tx : \star\right]; \left(\Gammabot, \tx : [= t]\right) \vdash \overline{B} : \overline{\mathcal{D}} }
        \infer1{ P; \Sigma; \Gammatop; \Gammabot \vdash \left(\opaque \tx = t, \overline{B}\right) : \left(\tx : \star, \overline{\mathcal{D}}\right) }
    \end{prooftree}
\] \\

\[
    \begin{prooftree}
        \hypo{ \Pi; \Sigma; \Gammatop\left[\Pi \texttt{+=} X : \Lambda (\overline{\tx : \star}). \{\}\right]; \overline{\tx : \star} \vdash \overline{B} : \overline{\mathcal{D}} }
        \hypo{ \Pi; \Sigma; \Gammatop\left[\Pi \texttt{+=} X : \Moddecl\right]; \Gammabot \vdash \overline{B'} : \overline{\mathcal{D}'} }
        \infer2{ \Pi; \Sigma; \Gammatop; \Gammabot \vdash \left(X : \Lambda (\overline{\tx : \star}). \{\overline{B}\}_{\overline{\BX[\overline{\tx'  = t}]}}, \overline{B'}\right) : \left(X : \Moddecl, \overline{\mathcal{D}'}\right) }
    \end{prooftree}
\] \\

\iffalse
    \section{Draft: capture-avoiding substitution}

    The substitution does not handle alpha-renaming itself.
    It assumes renaming is taken care of upstream.

    \begin{align*}
        \emptyenv[\tx := t]                  & = \emptyenv                                      & \\
        (\Gammabot, D)[\tx := t]             & = \Gammabot[\tx := t], D[\tx := t]               & \\
        (v : \bigcap \overline{T})[\tx := t] & = v[\tx := t] : (\bigcap \overline{T})[\tx := t] & \\
    \end{align*}

    \begin{align*}
        (\BX : \Lambda (\overline{\tx : \star}) \{\overline{D}\})[\tx'' := t] & = \{\overline{D}[\tx'' := t] \}                                                    & \\
        (X : \Moddecl)[\tx'' := t']                                           & = \{\overline{D}[\tx'' := t']\}_{\overline{\BX[\overline{\tx' = t[\tx'' := t']}]}} & \\
    \end{align*}

    \begin{align*}
        (X : t)[\tx := t']                    & = X : t[\tx := t']                    &                           \\
        (x : \bigcap \overline{T})[\tx := t'] & = x : \bigcap \overline{T[\tx := t']} &                           \\
        (\tx' : \star)[\tx := t']             & = \tx' : \star                        & \text{ if } \tx' \neq \tx \\
        (\tx' : [= t])[\tx := t']             & = \tx' : [= t[\tx := t']]             & \text{ if } \tx' \neq \tx \\
    \end{align*}


    \begin{align*}
        (\bigcap \overline{T})[\tx := t']                                     & =\bigcap \overline{T[\tx := t']}                                      &                                   \\
        ((\overline{N : T}) \rightarrow T')[\tx := t']                        & = ((\overline{N : T})[\tx := t']) \rightarrow T'[\tx := t']           &                                   \\
        (\{ \overline{D} \}_{\overline{\BX[\overline{\tx' = t}]}})[\tx := t'] & = \{ \overline{D} \}_{\overline{\BX[\overline{\tx' = t[\tx := t']}]}} &                                   \\
        \atom[\tx := t']                                                      & = \atom                                                               &                                   \\
        \mod[\tx := t']                                                       & = \mod                                                                &                                   \\
        \fun[\tx := t']                                                       & =\fun                                                                 &                                   \\
        ((\overline{t}) \rightarrow t_r)[\tx := t']                           & = (\overline{t[\tx := t']}) \rightarrow t_r[\tx := t']                &                                   \\
        \%\{\overline{k : t}\}[\tx := t']                                     & = \%\{\overline{k : t[\tx := t']}  \}                                 &                                   \\
        (t_1 \vee t_2)[\tx := t']                                             & = t_1[\tx := t'] \vee t_2[\tx := t']                                  &                                   \\
        (\lnot t)[\tx := t']                                                  & = \lnot (t[\tx := t'])                                                &                                   \\
        \alpha[\tx := t']                                                     & = \alpha                                                              &                                   \\
        \mathbb{O}[\tx := t']                                                 & = \mathbb{O}                                                          &                                   \\
        \tx[\tx := t']                                                        & = t'                                                                  &                                   \\
        \tx'[\tx := t']                                                       & = \tx'                                                                & \text{ if } \tx' \neq \tx         \\
        (M.\tx')[\tx := t']                                                   & = M.\tx'                                                              & \text { no matter if } \tx' = \tx \\
        (\BX[\overline{\tx' = t}])[\tx := t']                                 & = \BX[\overline{\tx' = t[\tx := t']}]                                 &                                   \\
        \langle M \rangle [\tx := t']                                         & = \langle M[\tx := t'] \rangle                                        &
    \end{align*}


    \begin{align*}
        X[\tx' := t']                               & = X                                         & \\
        (\Elixir.P[\overline{\tx = t}])[\tx' := t'] & = \Elixir.P[\overline{\tx = t[\tx' := t']}] & \\
    \end{align*}
\fi
\end{document}
